\begin{textsample}{NEG Dim 9 – global\_north\_2024\_05 – Score -2.00 – global\_north\_2024\_05/108446451.txt}  \label{ex:f9_neg_159}
The US and the UK will reject the international court of justice order directing Israel to end its offensive on Rafah after slowly blurring their red lines that once stated that they could not support a military offensive in Rafah.
The deputy foreign secretary, Andrew Mitchell, told MPs on Monday"the UK could only support a constructive plan for Rafah that complies with international humanitarian law on all counts".
On Tuesday he told the UK business select committee that"the significant operation in Rafah, it appears, has not yet started", even though 800,000 people had fled the area, including 400,000 who had been warned to do so by the Israel Defense Forces.
His definition of a major offensive -- which did not encompass an operation that led to the collective flight of so many people -- stretched the credulity of Labour MPs on the committee.
But he persisted, saying :"What we said is that we do not think an operation in Rafah should go ahead without there being a proper plan, and that we have seen no @ @ @ @ @ @ @ @ @ @ without seeing that plan, it should not go ahead."The select committee chair, Liam Byrne, cited the movement of 800,000 people :"If that is not significant, then what is?"Mitchell replied that the UK was doing what it could do to help with aid, adding the fact"800,000 had chosen to go of itself would not lead us to make a change in the assessment"of whether a serious breach of IHL had occurred.
The Labour MP Andy McDonald, previewing the line taken by the ICJ, asked the minister :"What choice did they have to move?
Was this just : ' I think I want to go and live somewhere else '?
Is that not a preposterous suggestion to make -- that this is a matter of free will?"Mitchell replied :"They have moved as a result of the circumstances."Byrne asked the minister directly :"Do you believe, Mr Mitchell, as the minister, that @ @ @ @ @ @ @ @ @ @ law in Rafah?"Mitchell replied :"It does not matter, chair, what I believe.
What is important is the legal process that informs that decision."He later admitted that the last released assessment of Israel ' s compliance rested on evidence that ended in January.
The US national security adviser, Jake Sullivan, took a different line in a briefing with reporters on Wednesday, indicating that he had been briefed by Israeli officials and by Israeli professionals on refinements to the Rafah plan that would achieve its military objectives while taking account of civilian harm.
Our morning email breaks down the key stories of the day, telling you what ' s happening and why it matters Enter your email address Privacy Notice : \textbf{Newsletters} may contain info about charities, online ads, and content funded by outside parties.
For more information see our Privacy Policy.
We use Google reCaptcha to protect our website and the Google Privacy Policy and Terms of Service apply. after \textbf{newsletter} promotion"@ @ @ @ @ @ @ @ @ @ ' s military operations in that area has been more targeted and limited, has not involved major military operations into the heart of dense urban areas,"he said."We now have to see what unfolds from here.
We will watch that, we will consider that, and we will see whether what Israel has briefed us and what they have laid out continues or if something else happens."He said there was no mathematical formula to decide if a plan was acceptable."What we ' re going to be looking at is whether there is a lot of death and destruction from this operation or if it is more precise and proportional,"he said.
He made no reference to the conditions in which Palestinians forced to flee were living.
It seems, according to interpretation, that the US either feels it has persuaded Israel to adjust its plans to make them acceptable or, faced with an Israeli fait accompli that the invasion would proceed regardless of Washington ' s objections, the US @ @ @ @ @ @ @ @ @ @ been that the threat to oppose a Rafah invasion was \textbf{useful} in trying to get both sides to agree to a ceasefire, but when those talks collapsed, the US administration saw no alternative to the Israeli offensive that removes what Israel regards as the last four Hamas battalions.
The US and the UK have also shifted ground on responsibility for the lack of aid reaching Gaza, another central part of the ICJ ruling.
Mitchell told MPs on Monday that the amount of aid reaching Gaza was"woefully inadequate".
The foreign secretary, David Cameron, said :"While there has been some progress in some areas of humanitarian relief, Israel must do more to make good its promises, and I am pressing them on this directly."But the US and the UK are now saying aid is not entering Rafah because of a dispute between Egypt and Israel, placing no greater responsibility on Israel even though it is the occupying power.

% matched lemmas: newsletter, useful
\end{textsample}
